% GENERATED FILE. Edit canonical lessons, then run npm run generate:books.
% canonical-source-hash: fnv1a64:8fc9eea00c2e4eb4
% canonical-lessons: HI-C496-anidra, HI-C496-uninda, HI-C496-pasarna, HI-C496-budbudana, HI-C496-bichauna

\chapter{Sleep}
\label{ch:hi-sleep}
\hlchaptermodality{496}

\begin{chapteropening}
\textbf{By the end of this chapter:} \emph{I can say how well someone slept.}

Complete the last lesson of chapter 496: I can say how well someone slept.
\end{chapteropening}

\section[\texorpdfstring{\hi{अनिद्रा} (anidrā)}{अनिद्रा (anidrā)}]{\hi{अनिद्रा} (anidrā) — insomnia; sleeplessness}

\label{lesson:HI-C496-anidra}



\begin{quote}
Before the new one: say the Hindi for nourishing, then the Hindi for to be digested.
\end{quote}

\subsection*{You'll want to know: \hi{अनिद्रा}}
\textbf{\hi{अनिद्रा}} — \emph{anidrā} — \textquotedblleft{}insomnia; sleeplessness\textquotedblright{}.

\begin{quote}
\hi{तनाव} \hi{से} \hi{अरुण} \hi{को} \hi{अनिद्रा} \hi{की} \hi{शिकायत} \hi{है।}
\end{quote}

\begin{quote}
\emph{tanāv se Aruṇ ko anidrā kī śikāyat hai.}
\end{quote}

\begin{quote}
Because of stress, Arun suffers from insomnia.
\end{quote}

Feminine.

\begin{cousinweb}
\textbf{\hi{अनिद्रा}} (\emph{anidrā}) is a tatsama, borrowed from Sanskrit \emph{a-nidrā}, \textquotedblleft{}sleeplessness\textquotedblright{}, built on \emph{a-}, \textquotedblleft{}not\textquotedblright{}, and \emph{nidrā}, \textquotedblleft{}sleep\textquotedblright{} (Monier-Williams).
\end{cousinweb}

\begin{grammarlens}[title={What you've built}]
One more word for describing what you lived and felt.
\end{grammarlens}

\subsection*{Your turn}
\begin{itemize}
\raggedright
  \item \emph{Say it:} \emph{anidrā}
  \item \emph{Say it:} \emph{anidrā}, once more
  \item \emph{Say it:} \emph{tanāv se Aruṇ ko anidrā kī śikāyat hai.} — Because of stress, Arun suffers from insomnia.
  \item \emph{Recall:} say the Hindi for nourishing, then the Hindi for to be digested, then say \emph{anidrā} again
\end{itemize}

\begin{tcolorbox}[breakable,colback=teal!4,colframe=teal!35!black,title={Before you move on}]
What does \hi{अनिद्रा} mean? (\textquotedblleft{}Insomnia; sleeplessness\textquotedblright{}.) Say it once more.
\end{tcolorbox}
\section[\texorpdfstring{\hi{उनींदा} (unīndā)}{उनींदा (unīndā)}]{\hi{उनींदा} (unīndā) — drowsy; heavy with sleep}

\label{lesson:HI-C496-uninda}



\begin{quote}
Before the new one: say the Hindi for to be digested, then the Hindi for insomnia.
\end{quote}

\subsection*{You'll want to know: \hi{उनींदा}}
\textbf{\hi{उनींदा}} — \emph{unīndā} — \textquotedblleft{}drowsy; heavy with sleep\textquotedblright{}.

\begin{quote}
\hi{सुबह} \hi{छह} \hi{बजे} \hi{बच्चे} \hi{उनींदे} \hi{थे।}
\end{quote}

\begin{quote}
\emph{subah chah baje bacce unīnde the.}
\end{quote}

\begin{quote}
At six in the morning the children were still drowsy.
\end{quote}

\begin{cousinweb}
\textbf{\hi{उनींदा}} (\emph{unīndā}) is inherited: CDIAL lists Hindi \emph{unī̃dā}, with the sense \textquotedblleft{}sleepless, with eyes open\textquotedblright{}, under Sanskrit \emph{unnidra}, \textquotedblleft{}sleepless\textquotedblright{} (CDIAL 2108). Today it means eyes heavy from too little sleep, a turn CDIAL does not give.
\end{cousinweb}

\begin{grammarlens}[title={What you've built}]
One more word for describing what you lived and felt.
\end{grammarlens}

\subsection*{Your turn}
\begin{itemize}
\raggedright
  \item \emph{Say it:} \emph{unīndā}
  \item \emph{Say it:} \emph{unīndā}, once more
  \item \emph{Say it:} \emph{subah chah baje bacce unīnde the.} — At six in the morning the children were still drowsy.
  \item \emph{Recall:} say the Hindi for to be digested, then the Hindi for insomnia, then say \emph{unīndā} again
\end{itemize}

\begin{tcolorbox}[breakable,colback=teal!4,colframe=teal!35!black,title={Before you move on}]
What does \hi{उनींदा} mean? (\textquotedblleft{}Drowsy; heavy with sleep\textquotedblright{}.) Say it once more.
\end{tcolorbox}
\section[\texorpdfstring{\hi{पसरना} (pasarnā)}{पसरना (pasarnā)}]{\hi{पसरना} (pasarnā) — to sprawl; to stretch out}

\label{lesson:HI-C496-pasarna}



\begin{quote}
Before the new one: say the Hindi for insomnia, then the Hindi for drowsy.
\end{quote}

\subsection*{You'll want to know: \hi{पसरना}}
\textbf{\hi{पसरना}} — \emph{pasarnā} — \textquotedblleft{}to sprawl; to stretch out\textquotedblright{}.

\begin{quote}
\hi{काम} \hi{से} \hi{लौटकर} \hi{अरुण} \hi{बिस्तर} \hi{पर} \hi{पसर} \hi{गया।}
\end{quote}

\begin{quote}
\emph{kām se lauṭkar Aruṇ bistar par pasar gayā.}
\end{quote}

\begin{quote}
Back from work, Arun sprawled out on the bed.
\end{quote}

Past (he …): \hi{वह} \hi{बिस्तर} \hi{पर} \hi{पसर} \hi{गया}, \emph{vah bistar par pasar gayā}.

\begin{cousinweb}
\textbf{\hi{पसरना}} (\emph{pasarnā}) is a tadbhava: CDIAL lists Hindi \emph{pasarnā}, \textquotedblleft{}to spread out\textquotedblright{}, under Sanskrit \emph{prásarati}, \textquotedblleft{}spreads out\textquotedblright{} (CDIAL 8825): a tired body spreads out across the bed.
\end{cousinweb}

\begin{grammarlens}[title={What you've built}]
One more word for describing what you lived and felt.
\end{grammarlens}

\subsection*{Your turn}
\begin{itemize}
\raggedright
  \item \emph{Say it:} \emph{pasarnā}
  \item \emph{Say it:} \emph{pasarnā}, once more
  \item \emph{Say it:} \emph{kām se lauṭkar Aruṇ bistar par pasar gayā.} — Back from work, Arun sprawled out on the bed.
  \item \emph{Recall:} say the Hindi for insomnia, then the Hindi for drowsy, then say \emph{pasarnā} again
\end{itemize}

\begin{tcolorbox}[breakable,colback=teal!4,colframe=teal!35!black,title={Before you move on}]
What does \hi{पसरना} mean? (\textquotedblleft{}To sprawl; to stretch out\textquotedblright{}.) Say it once more.
\end{tcolorbox}
\section[\texorpdfstring{\hi{बुदबुदाना} (budbudānā)}{बुदबुदाना (budbudānā)}]{\hi{बुदबुदाना} (budbudānā) — to mumble}

\label{lesson:HI-C496-budbudana}



\begin{quote}
Before the new one: say the Hindi for drowsy, then the Hindi for to sprawl.
\end{quote}

\subsection*{You'll want to know: \hi{बुदबुदाना}}
\textbf{\hi{बुदबुदाना}} — \emph{budbudānā} — \textquotedblleft{}to mumble\textquotedblright{}.

\begin{quote}
\hi{बुख़ार} \hi{में} \hi{दादा} \hi{कुछ} \hi{बुदबुदा} \hi{रहे} \hi{थे।}
\end{quote}

\begin{quote}
\emph{bukhār meṁ dādā kuch budbudā rahe the.}
\end{quote}

\begin{quote}
In his fever my grandfather was mumbling something.
\end{quote}

Past (he …): \hi{वह} \hi{कुछ} \hi{बुदबुदाया}, \emph{vah kuch budbudāyā}.

\begin{cousinweb}
\textbf{\hi{बुदबुदाना}} (\emph{budbudānā}) is a tadbhava: CDIAL lists Hindi \emph{budbudānā}, \textquotedblleft{}to mutter\textquotedblright{}, under \emph{*budabuda}, a form it reconstructs (CDIAL 9274); the doubled syllable is the sound of lips moving with no clear words.
\end{cousinweb}

\begin{grammarlens}[title={What you've built}]
One more word for describing what you lived and felt.
\end{grammarlens}

\subsection*{Your turn}
\begin{itemize}
\raggedright
  \item \emph{Say it:} \emph{budbudānā}
  \item \emph{Say it:} \emph{budbudānā}, once more
  \item \emph{Say it:} \emph{bukhār meṁ dādā kuch budbudā rahe the.} — In his fever my grandfather was mumbling something.
  \item \emph{Recall:} say the Hindi for drowsy, then the Hindi for to sprawl, then say \emph{budbudānā} again
\end{itemize}

\begin{tcolorbox}[breakable,colback=teal!4,colframe=teal!35!black,title={Before you move on}]
What does \hi{बुदबुदाना} mean? (\textquotedblleft{}To mumble\textquotedblright{}.) Say it once more.
\end{tcolorbox}
\section[\texorpdfstring{\hi{बिछौना} (bichaunā)}{बिछौना (bichaunā)}]{\hi{बिछौना} (bichaunā) — a bedroll; bedding laid out}

\label{lesson:HI-C496-bichauna}



\begin{quote}
Before the new one: say the Hindi for to sprawl, then the Hindi for to mumble.
\end{quote}

\subsection*{You'll want to know: \hi{बिछौना}}
\textbf{\hi{बिछौना}} — \emph{bichaunā} — \textquotedblleft{}a bedroll; bedding laid out\textquotedblright{}.

\begin{quote}
\hi{दादी} \hi{ने} \hi{मेहमानों} \hi{के} \hi{लिए} \hi{फ़र्श} \hi{पर} \hi{बिछौना} \hi{बिछाया।}
\end{quote}

\begin{quote}
\emph{dādī ne mehmānoṁ ke liye farś par bichaunā bichāyā.}
\end{quote}

\begin{quote}
Grandmother laid out bedding on the floor for the guests.
\end{quote}

Masculine.

\begin{cousinweb}
\textbf{\hi{बिछौना}} (\emph{bichaunā}) is a tadbhava: CDIAL lists Hindi \emph{bichaunā}, \textquotedblleft{}bedding\textquotedblright{}, under \emph{*vicchādana}, a form it reconstructs (CDIAL 11691): what is spread out to sleep on.
\end{cousinweb}

\begin{grammarlens}[title={What you've built}]
One more word for describing what you lived and felt.
\end{grammarlens}

\subsection*{Your turn}
\begin{itemize}
\raggedright
  \item \emph{Say it:} \emph{bichaunā}
  \item \emph{Say it:} \emph{bichaunā}, once more
  \item \emph{Say it:} \emph{dādī ne mehmānoṁ ke liye farś par bichaunā bichāyā.} — Grandmother laid out bedding on the floor for the guests.
  \item \emph{Recall:} say the Hindi for to sprawl, then the Hindi for to mumble, then say \emph{bichaunā} again
\end{itemize}

\begin{tcolorbox}[breakable,colback=teal!4,colframe=teal!35!black,title={Before you move on}]
What does \hi{बिछौना} mean? (\textquotedblleft{}A bedroll; bedding laid out\textquotedblright{}.) Say it once more.
\end{tcolorbox}
